% Introduction to Part III. Unnumbered text: no \label except the \phantomsection target.
\phantomsection\label{intro:III}
Part~III replaces the global hypothesis~(B) of Part~II by Ba\l aban's local one. A real background is $(M,\eps)$-regular if on every cube of side $8M+4$ there is a gauge in which $U=e^{\ii\eta A}$ with $|A|$, $|\partial A|$ and the H\"older quotients of $\partial A$ at most $\eps$; the gauges on different cubes are unrelated (Definition~\ref{def:regular}). Part~VI shows that the background fields, and all real configurations with small curvature and current, are regular (Theorem~\ref{thm:balaban}, Corollary~\ref{cor:regularity}). The variational problem of Part~V and the step of Part~VII work at such backgrounds, and Parts~VIII and~IX build their localized operators from the Dirichlet operators of this Part.

\paragraph{What this Part proves.} The main result is Theorem~\ref{thm:G}. For $M\ge M_0$ and $\eps$ small, independently of $M$:
\begin{itemize}
\item the operators of Part~II exist at regular backgrounds and keep their bounds, without the comparisons with $U=1$, with covariant derivatives and with H\"older quotients in the local gauges (Theorems~\ref{thm:reg-scalar}, \ref{thm:reg-gauge} and~\ref{thm:reg-min});
\item the scalar and the gauge-field operator have Dirichlet versions on unions $X$ of blocks, which depend only on the background on the bonds of $\hat X$, the points at $\infty$-distance at most $1$ from $X$, and are coercive; their inverses, the inverse of the averaged scalar Green function and $\one_X-R_X$ decay exponentially, uniformly in the domain (Theorems~\ref{thm:reg-scalar}, \ref{thm:reg-coercive} and~\ref{thm:reg-gauge}); the gauge-field version carries a mass term near the boundary and Dirichlet conditions for all components;
\item if two regular backgrounds agree on a region, the operators agree deep inside it, up to errors exponentially small in the distance to the complement (Theorem~\ref{thm:reg-locality});
\item all of this extends analytically to small complex perturbations $Ue^{\ii\eta\,\delta\!A}$, with $\Rea$-coercivity in place of coercivity (Proposition~\ref{prop:reg-complex}), and, for $M$ large, Ba\l aban's random-walk expansion converges for the scalar Green function (Theorem~\ref{thm:random-walk}(a)).
\end{itemize}
The key step is a local proof of the coercivity of the gauge-field operator (Theorem~\ref{thm:reg-coercive}): an IMS localization at scale $M$ (Lemma~\ref{lem:IMS}) reduces it to cubes, where one works in the local gauge and uses the coercivity of Part~I. The statements are proved for the model pair of Part~II; Remark~\ref{rem:general-averaging} lists the properties of an averaging pair that the proofs use, and Corollary~\ref{cor:parts-apply} verifies them for Ba\l aban's averaging.

\paragraph{Imports.} Part~II throughout, and from Part~I the coercivity (Corollary~\ref{cor:scalar-coercive}, Theorem~\ref{thm:gauge-coercive}), the Combes--Thomas lemmas and Theorem~\ref{thm:local-regularity}. Four passages point ahead. The proof of Lemma~\ref{lem:full-gradient} uses the covariant Hodge inequality of Lemma~\ref{lem:weitzenbock}, a summation by parts that uses only the curvature bound. Theorem~\ref{thm:random-walk}(b), the walk for the gauge-field Green function, is Theorem~\ref{thm:weak-inverse}(a) at level~(L3) with $s=1$; it is proved with Section~\ref{sec:nested-walks} and is best read after it. The second paragraph of Proposition~\ref{prop:reg-complex}, and its proof, concern the potentials of Section~\ref{sec:dirichlet-averaged} (Part~VIII). The last paragraph of Remark~\ref{rem:general-averaging} uses Propositions~\ref{prop:AB} and~\ref{prop:complex-ext} of Part~X.

\paragraph{Exports.} Definition~\ref{def:regular} and Theorem~\ref{thm:reg-scalar} are used in Parts~IV--XI, and Theorems~\ref{thm:reg-coercive}, \ref{thm:reg-gauge} and~\ref{thm:reg-min} in Parts~IV--IX. Locality and the comparison of domains (Theorem~\ref{thm:reg-locality}, Lemma~\ref{lem:compare-inverse}) are the base of the localization in Parts~VII--IX, and Proposition~\ref{prop:reg-complex} is used in Parts~IV and VIII--X. Lemma~\ref{lem:uniformity}(a) shows that the constants of Theorems~\ref{thm:reg-scalar}, \ref{thm:reg-coercive} and~\ref{thm:reg-gauge} (on the torus) and~\ref{thm:reg-min} do not depend on $M$, $L$, $k$ or $n$; the random-walk expansion and Theorem~\ref{thm:reg-locality}(b) have constants that depend on $M$ (Remark~\ref{rem:uniformity-limits}(b)). The bounds of Parts~I--VIII do not use the random-walk expansion (Remark~\ref{x:20.3}); the weakened walks of Part~IX are built on the same kind of parametrix (Lemma~\ref{lem:parametrix-X}), and there the walk series also gives the bounds (Theorem~\ref{thm:weak-inverse}).

\paragraph{Details for this Part.} The companion has no chapter on Part~III; its parameter ledger lists the constants of Part~III that Parts~VII--XI use, among them $M_0$.
